\chapter{Problem Statement and Requirements}
\label{chap:problem}

This chapter turns the gap established in Chapter~\ref{chap:background}
into a concrete engineering problem: a target scenario
(\S\ref{sec:problem-scenario}), a set of requirements that any solution
must satisfy (\S\ref{sec:problem-requirements}), and the boundaries of
what this thesis does not attempt (\S\ref{sec:problem-nongoals}). The
requirements are numbered R1--R7 and are used as the evaluation rubric in
Chapters~\ref{chap:evaluation} and~\ref{chap:conclusions}.

\section{Target Scenario}
\label{sec:problem-scenario}

The setting is a multi-agent LLM workload on an HPC cluster. A user holds
a SLURM allocation of $M$ compute nodes (in this thesis's experiments,
$M \in \{4, 8\}$ on the ARES cluster, with analyses exercised up to
several hundred simulated nodes). On each node run $K$ LLM agents:
independent processes, each an agent loop over its own conversation
context, each identified by a scenario-scoped session identifier of the
form \texttt{role@scenario}. Agents cooperate: an agent may delegate a
sub-task to a peer on another node by invoking a remote-agent MCP tool,
producing nested chains of cross-host calls. A ChronoLog deployment
shares the allocation: a keeper on every node, a visor coordinating, and
a grapher draining sealed chunks to an archive on shared storage.

The human in the loop is the \emph{operator}, the researcher or
engineer running the fleet, who needs answers both while the system runs
and after it finishes:

\begin{enumerate}
\item[Q1.] \emph{What is the swarm doing right now?} Which agents exist,
  which are talking to which, and on which physical nodes?
\item[Q2.] \emph{Is anything failing?} Not as a scroll of per-event
  errors, but as a ranked, deduplicated set of problems with causes.
\item[Q3.] \emph{Why is this scenario slow?} Which hop in a cross-host
  delegation chain is the bottleneck?
\item[Q4.] \emph{Which node is dragging the fleet down?} A
  population-level comparison across all $M$ nodes, including silent
  ones.
\item[Q5.] \emph{Has this happened before?} Whether a failure is novel
  or a recurring pattern across runs.
\end{enumerate}

Q1--Q3 are agent-semantic questions; Q4 is a fleet-topology question; Q5
is a longitudinal question. Q2--Q4 in combination require pivoting
between layers (from a fleet symptom to a single agent's
conversation), the pivot neither tooling camp of
\S\ref{sec:bg-sota} supports.

\section{Requirements}
\label{sec:problem-requirements}

\paragraph{R1: Dual-plane causal capture.} The system must capture
both observable planes of \S\ref{sec:bg-agent-signals}: LLM turns
(prompt, response, model, tokens, cost, latency, working-memory deltas)
and inter-agent calls (caller, callee, hosts, timing, outcome). The two
planes must also be \emph{joinable}: an inter-agent edge must resolve
to the conversations of its endpoint agents, and every event must be
attributed to the physical node that produced it. Without the join, Q1
and the Q2--Q4 pivot are unanswerable.

\paragraph{R2: Real-time visibility.} Live views (Q1, and the live
half of Q2/Q4) must reflect events within seconds of ingest. In
particular, the system must not be limited by the substrate's drain
latency (\S\ref{sec:bg-chronolog-design}), and no read may ever block
the serving process on un-drained log state.

\paragraph{R3: Durability and replayability.} Every captured event
must survive process and node restarts and remain replayable after the
run ends: post-hoc analysis, cost accounting, and the longitudinal
question Q5 all read history. Transient failures of downstream
components must not lose events (at-least-once capture), and retries
must not duplicate them in the durable store (effectively exactly-once
ingest, \S\ref{sec:bg-delivery}).

\paragraph{R4: Fleet scale.} Population-level views (Q4) must remain
responsive at 100+ nodes. Concretely, the cost of rendering the fleet
overview must not grow with the number of captured events, only with the
number of nodes and the length of the queried window. Silent
nodes must also appear, since silence is itself a signal.

\paragraph{R5: Minimal, non-blocking overhead.} Instrumentation must
not perturb the workload it observes (the Dapper tenet,
\S\ref{sec:bg-tracing-lineage}). The agent-side fast path must be a
local, constant-time operation (an append to node-local state),
with all network and log I/O off the agent's critical path. The
capture machinery must degrade gracefully (buffer, fall back) and never
back-pressure an agent.

\paragraph{R6: Zero-marginal-cost diagnosis.} Failure detection and
diagnosis (Q2, Q5) must not itself consume LLM tokens by default:
a diagnosis pipeline that costs money per scan will be disabled in
practice. LLM-assisted diagnosis may exist, but only as an opt-in on
top of a deterministic baseline.

\paragraph{R7: Drop-in deployment.} The system must attach to an
existing ChronoLog deployment without modifying ChronoLog, the
scheduler, or the agent framework: instrumentation via a thin client
library or local HTTP endpoint, per-node components launchable on a
plain SLURM allocation, and a read-only dashboard that can point at any
visor.

\section{Non-Goals}
\label{sec:problem-nongoals}

Four adjacent problems are explicitly out of scope. The system is not
an \emph{agent framework}: it observes agents built with existing SDKs
and protocols, and takes no part in orchestrating them. It is not a
\emph{scheduler or resource manager}: placement, allocation, and job
control remain SLURM's, and the system only queries it. And it is not a
general \emph{metrics time-series database}: it captures the
agent-semantic event planes of R1, not arbitrary node telemetry.
Where classical infrastructure metrics are needed, the systems of
\S\ref{sec:bg-sota-hpc} remain the right tool, and the two can run side
by side. Finally, the thesis targets observation and diagnosis, not
automated remediation: acting on a diagnosis (restarting agents, rolling
back context) is surfaced as a recommendation to the operator, not
executed autonomously.

With the questions, requirements, and non-goals fixed, the problem is
fully specified. Chapter~\ref{chap:design} presents a design that meets
R1--R7 on the ChronoLog substrate; Chapter~\ref{chap:evaluation} then
measures each requirement against it.
